\documentclass[11pt]{article}
\usepackage[T1]{fontenc}
\usepackage{lmodern}
\usepackage[margin=27mm]{geometry}
\usepackage{amsmath,amssymb,amsthm}
\usepackage{booktabs,longtable,array}
\usepackage{listings}
\usepackage[obeyspaces,spaces]{url}
\usepackage{microtype}
\usepackage[hidelinks]{hyperref}
% The author's macros.
%
% some calligraphy
%
\newcommand\cA{{\mathcal A}}
\newcommand\cB{{\mathcal B}}
\newcommand\cC{{\mathcal C}}
\newcommand\cD{{\mathcal D}}
\newcommand\cE{{\mathcal E}}
\newcommand\cF{{\mathcal F}}
\newcommand\cG{{\mathcal G}}
\newcommand\cH{{\mathcal H}}
\newcommand\cI{{\mathcal I}}
\newcommand\cJ{{\mathcal J}}
\newcommand\cK{{\mathcal K}}
\newcommand\cL{{\mathcal L}}
\newcommand\cM{{\mathcal M}}
\newcommand\cN{{\mathcal N}}
\newcommand\cP{{\mathcal P}}
\newcommand\cQ{{\mathcal Q}}
\newcommand\cR{{\mathcal R}}
\newcommand\cS{{\mathcal S}}
\newcommand\cT{{\mathcal T}}
\newcommand\cU{{\mathcal U}}
\newcommand\cV{{\mathcal V}}

%
% fields and rings (and a semigroup)
%
\newcommand\cx{{\mathbb C}}% complexes
\newcommand\fld{{\mathbb F}}
\newcommand\flde{{\mathbb E}}
\newcommand\ints{{\mathbb Z}}
\newcommand\nn{{\mathbb N}}%non-negative integers
\newcommand\re{{\mathbb R}}%reals
\newcommand\rats{{\mathbb Q}}

\newcommand\Zu{{\mathbf u}}
\newcommand\Zv{{\mathbf v}}
\newcommand\Zw{{\mathbf w}}
\newcommand\Zy{{\mathbf y}}
\newcommand\Zz{{\mathbf z}}
\newcommand\Zx{{\mathbf x}}
\newcommand\Zj{{\mathbf j}}


%
% the really useful stuff
%
%\newcommand\opk[1]{\mathop{\hbox{\rm #1}}\nolimits} %% old version
\newcommand\opk[1]{\mathop{\mathrm{#1}}\nolimits}
\newcommand\comp[1]{{\mkern2mu\overline{\mkern-2mu#1}}}
\newcommand\diff{\mathbin{\mkern-1.5mu\setminus\mkern-1.5mu}}% for \setminus
\newcommand\gbin[2]{\genfrac{[}{]}{0pt}{}{#1}{#2}}
\newcommand\inv{^{-1}}
\newcommand\res{\mathbin{\mkern-2.0mu\restriction\mkern-2.0mu}}% ch2
\newcommand\sbs{\subseteq}
\newcommand\sps{\supseteq}
\newcommand\seq[3]{#1_{#2},\ldots,#1_{#3}}
\newcommand\supp{\opk{supp}}
\newcommand\im{\opk{im}}
\newcommand\row{\opk{row}}
\newcommand\pmat[1]{\begin{pmatrix} #1 \end{pmatrix}}
\newcommand\cprod{\mathbin{\square}}
%
% matrix theory
%
\newcommand\ip[2]{\langle#1,#2\rangle}
\newcommand\one{{\bf1}}
\newcommand\rk{\opk{rk}}
\newcommand\tr{\opk{tr}}
\newcommand\col{\opk{col}}
\newcommand\mat[3]{\mathrm{Mat}_{#1\times #2}(#3)}

%
% some group theory
%
\newcommand\aut[1]{{\rm Aut}(#1)}
\newcommand\fx[1]{{\rm fix}(#1)}% ch2
\newcommand\grp[1]{\langle #1\rangle}
\newcommand\nrml{\vartriangleleft}
\newcommand\nrmleq{\trianglelefteq}
\newcommand\alt[1]{\opk{Alt}(#1)}
\DeclareMathOperator{\Sym}{Sym}
% End of the author's macros.
\newtheorem{theorem}{Theorem}[section]
\newtheorem{lemma}[theorem]{Lemma}
\newtheorem{corollary}[theorem]{Corollary}
\theoremstyle{definition}
\newtheorem{definition}[theorem]{Definition}
\newtheorem{remark}[theorem]{Remark}
\newcommand{\Ze}{{\mathbf e}}
\newcommand{\code}[1]{{\small\nolinkurl{#1}}}
\newcolumntype{L}[1]{>{\raggedright\arraybackslash}p{#1}}
\lstset{basicstyle=\ttfamily\small,columns=fullflexible,
  keepspaces=true,breaklines=true,showstringspaces=false}
\title{There is no $2$-$(15,5,2)$ design}
\author{}
\date{28 September 2026}

\begin{document}
\maketitle
\begin{abstract}
We prove that no $2$-$(15,5,2)$ design exists, allowing labelled repeated
blocks. We follow Hall and Connor's embedding argument at block size five
and prove each result used from Connor. The hypothetical design extends
to a symmetric $2$-$(22,7,2)$ design. Its incidence matrix would have an
integer determinant whose square is $49\cdot5^{21}$, a contradiction.
We give the finite matrix certificates, their checks, and the
Lean statements.
\end{abstract}

\section{The statement and conventions}

A \emph{$2$-$(v,k,\lambda)$ design} consists of a finite point set $P$, a
finite set $\cB$ of block labels, and a subset $B\sbs P$ for each label
$B\in\cB$. We use the same letter for a label and its subset. We require
$|P|=v$, $|B|=k$ for every block, and exactly $\lambda$ block labels
containing each pair of distinct points. The definition allows different
labels to have the same subset; Lemma~\ref{lem:intersections} will exclude
this for the parameters considered here. We will write $b=|\cB|$ and, when it is constant,
$r$ for the number of blocks containing a point. A design is
\emph{symmetric} when $b=v$.

\begin{theorem}\label{thm:main}
There is no $2$-$(15,5,2)$ design. Equivalently, there is no design with
$(v,k,b,r,\lambda)=(15,5,21,7,2)$.
\end{theorem}

This is the nonexistence result used in Remark 5.12 and Example 5.13 of
Cameron, Goethals, Seidel and Shult~\cite{CGSS}. They cite Hall's book;
we will use Hall and Connor~\cite{HC}, with all the necessary arguments
included here. Their ``derived design'' is a \emph{residual design} in
current terminology: we delete one block and all its points. We will
use the current terminology.

We will work over $\rats$ for the matrix arguments, keeping the incidence
matrices integral for the final determinant. We will use functional
notation for entries: $M(u,w)$ and $\Zx(u)$. We write $I_U$ for the
identity matrix on $U$, $J_{U,V}$ for the all-ones matrix with row set $U$
and column set $V$, and $J_U=J_{U,U}$. We omit index sets when their sizes
are determined by multiplication. Inner products are written
$(\Zx,\Zy)=\Zx^T\Zy$. All ranks and dimensions are over $\rats$.

Our proof follows HC1--HC6 of \code{cgss/prerequisites.md}. We will first
determine block intersections and construct Connor's matrix. Its
nonnegative quadratic form and rank bound give the local clique
structure. The family of all six-cliques then supplies the seven new
points of the embedding. Finally, we will prove the even-order
Bruck--Ryser--Chowla obstruction in the single instance needed here.

\section{Counting intersections: HC1 and HC2}

\begin{lemma}\label{lem:parameters}
Suppose $D$ is a $2$-$(15,5,2)$ design on $P$ with block labels $\cB$.
Every point belongs to seven blocks, and $|\cB|=21$.
\end{lemma}
\begin{proof}
Fix $p\in P$, and let $r_p$ be the number of blocks containing $p$.
Count pairs $(q,B)$ with $q\ne p$ and $p,q\in B$. Counting by $B$ gives
$4r_p$; counting by $q$ gives $14\cdot2$. Thus $r_p=7$. Counting
incident pairs $(p,B)$ now gives $5|\cB|=15\cdot7=105$, so
$|\cB|=21$. These counts concern block labels, so they also apply when
blocks repeat.
\end{proof}


The incidence matrix is $N\in\mat{P}{\cB}{\ints}$, with $N(p,B)=1$
if $p\in B$ and $0$ otherwise. Put $G=N^TN$; then
$G(B,B')=|B\cap B'|$.

\begin{lemma}\label{lem:matrix-identities}
Suppose $D$ is a $2$-$(15,5,2)$ design, with incidence matrix $N$ and
$G=N^TN$. Then
\begin{equation}\label{eq:incidence}
 NN^T=5I_P+2J_P,\qquad
 N\Zj_{\cB}=7\Zj_P,\qquad N^T\Zj_P=5\Zj_{\cB},
\end{equation}
and
\begin{equation}\label{eq:G-identities}
 G^2=5G+50J_{\cB},\qquad
 GJ_{\cB}=J_{\cB}G=35J_{\cB},\qquad J_{\cB}^2=21J_{\cB}.
\end{equation}
\end{lemma}
\begin{proof}
The diagonal entries of $NN^T$ count the seven blocks on a point; its
off-diagonal entries count the two blocks on a pair of points.
The other two identities in~\eqref{eq:incidence} are the row and column
sums of $N$. Thus
\[
 G^2=N^T(NN^T)N=5N^TN+2N^TJ_PN=5G+50J_{\cB},
\]
since every entry of $N^TJ_PN$ is $5\cdot5=25$.
Also $G\Zj_{\cB}=N^T(7\Zj_P)=35\Zj_{\cB}$. Symmetry gives the
other product with $J_{\cB}$. Finally, each entry of $J_{\cB}^2$ is
$|\cB|=21$.
\end{proof}

\begin{lemma}\label{lem:intersections}
\textup{\cite[Lemma 4.1, at $k=5$]{HC}}
Suppose $D$ is a $2$-$(15,5,2)$ design. Distinct block labels have
intersection size one or two. Every block has ten other blocks meeting
it in one point, and ten meeting it in two points.
\end{lemma}
\begin{proof}
Fix $B$, and for $0\le i\le5$ let $n_i$ count the other labels $B'$
with $|B\cap B'|=i$. Counting other labels, incident points of $B$, and
ordered pairs of distinct points of $B$, respectively, gives
\[
 \sum_{i=0}^5 n_i=20,\qquad
 \sum_{i=0}^5 i n_i=5(7-1)=30,\qquad
 \sum_{i=0}^5 i(i-1)n_i=5\cdot4(2-1)=20.
\]
For the last equality, every ordered pair of distinct points of $B$
lies in exactly one further block label. Taking this linear combination
of the three counts, in $\ints$, gives
\[
 \sum_{i=0}^5(i-1)(i-2)n_i=20-2\cdot30+2\cdot20=0.
\]
The coefficients at $i=0,1,2,3,4,5$ are $2,0,0,2,6,12$.
All summands are nonnegative, so $n_i=0$ unless $i=1$ or $i=2$.
The first two counts give $n_1+n_2=20$ and $n_1+2n_2=30$, whence
$n_1=n_2=10$. In particular, distinct block labels cannot represent
the same subset, since that would give intersection size five.
\end{proof}

We will denote by $Z$ the graph on $\cB$ in which distinct blocks are
adjacent when their intersection has size one. We write $B\sim B'$ for
adjacency and $\Gamma(B)$ for the neighbourhood of $B$. A \emph{clique}
is a set of vertices any two distinct members of which are adjacent.

\begin{lemma}\label{lem:common}
\textup{\cite[Lemma 4.2, at $k=5$]{HC}}
Suppose $D$ is a $2$-$(15,5,2)$ design and $B_1\ne B_2$ are blocks.
For $i,j\in\{1,2\}$ let $x_{ij}$ count the blocks $B\notin\{B_1,B_2\}$
with $|B_1\cap B|=i$ and $|B_2\cap B|=j$. The counts are
\[
\begin{array}{c|rrrr}
 |B_1\cap B_2|&x_{11}&x_{12}&x_{21}&x_{22}\\ \hline
 1&5&4&4&6\\
 2&4&6&6&3
\end{array}.
\]
In particular, $Z$ has degree ten, and adjacent vertices have five
common neighbours.
\end{lemma}
\begin{proof}
Put $s=|B_1\cap B_2|$. There are $19$ remaining block labels.
The sum of their intersections with either initial block is $30-s$,
by the preceding proof. By~\eqref{eq:G-identities}, the sum of the
products of these two intersections is
\[
 (G^2)(B_1,B_2)-5s-5s=(5s+50)-10s=50-5s.
\]
Partitioning the remaining labels by their two intersection sizes gives
\[
\begin{aligned}
 x_{11}+x_{12}+x_{21}+x_{22}&=19,\\
 x_{11}+x_{12}+2x_{21}+2x_{22}&=30-s,\\
 x_{11}+2x_{12}+x_{21}+2x_{22}&=30-s,\\
 x_{11}+2x_{12}+2x_{21}+4x_{22}&=50-5s.
\end{aligned}
\]
Thus
\[
\begin{aligned}
 x_{22}&=(50-5s)-2(30-s)+19=9-3s,\\
 x_{12}=x_{21}&=(30-s)-19-x_{22}=2+2s,\\
 x_{11}&=19-x_{12}-x_{21}-x_{22}=6-s.
\end{aligned}
\]
For example, the first identity sums $(i-1)(j-1)$ over the remaining
blocks, and the second sums $j-1$ or $i-1$. Substituting $s=1$ and
$s=2$ proves the table. The degree assertion is
Lemma~\ref{lem:intersections}.
\end{proof}

Thus $Z$ is strongly regular: it has $21$ vertices, valency ten,
five common neighbours for each adjacent pair and four for each
nonadjacent pair. We will reserve $k$ for the design's block size.
We will not use a classification of graphs with these parameters.

\section{Connor's matrix: HC3}

\begin{lemma}\label{lem:Connor}
Suppose $D$ is a $2$-$(15,5,2)$ design with incidence matrix $N$ and
$G=N^TN$. Define
\begin{equation}\label{eq:C-definition}
 C=35I_{\cB}-7G+10J_{\cB}.
\end{equation}
Then $C$ is symmetric,
\begin{equation}\label{eq:C-identities}
 C^2=35C,\qquad CN^T=0,
\end{equation}
and its entries are
\begin{equation}\label{eq:C-entries}
 C(B,B')=
 \begin{cases}
 10,&B=B';\\
 3,&B\sim B';\\
 -4,&B\ne B'\text{ and }B\not\sim B'.
 \end{cases}
\end{equation}
For every $\Zx\in\rats^{\cB}$,
\begin{equation}\label{eq:C-positive}
 35\Zx^TC\Zx=(C\Zx)^T(C\Zx)\ge0.
\end{equation}
Moreover, $\rk C\le6$.
\end{lemma}
\begin{proof}
Symmetry and the entries follow from the definition and
Lemma~\ref{lem:intersections}. Using~\eqref{eq:G-identities}, we obtain
\[
\begin{aligned}
 C^2
 &=1225I+49G^2+100J^2-490G+700J-140GJ\\
 &=1225I+49(5G+50J)+2100J-490G+700J-4900J\\
 &=1225I-245G+350J=35C.
\end{aligned}
\]
Here both $G$ and $J$ are indexed by $\cB$.
Also $N^TJ_P=5J_{\cB,P}$ and $J_{\cB}N^T=7J_{\cB,P}$, so
\[
\begin{aligned}
 CN^T
 &=35N^T-7N^T(NN^T)+10J_{\cB}N^T\\
 &=35N^T-7N^T(5I_P+2J_P)+70J_{\cB,P}=0.
\end{aligned}
\]
The first identity of~\eqref{eq:C-identities} and symmetry give
\eqref{eq:C-positive}; its right side is a sum of rational squares.

The determinant formula of Lemma~\ref{lem:det-rank-one}, whose proof
is independent of the design argument, gives
\[
 \det(NN^T)=\det(5I_{15}+2J_{15})=5^{14}\cdot35\ne0.
\]
Thus $NN^T$ is invertible and has rank $15$. Over $\rats$, the
matrix rank identities give $\rk N=\rk(NN^T)=15$ and
$\rk N^T=\rk N$. Since $CN^T=0$, the rank inequality for a zero
product gives
\[
 \rk C+\rk N^T\le|\cB|=21.
\]
Thus $\rk C\le21-15=6$. The same determinant formula will serve
again at $n=22$ in the final obstruction.
\end{proof}

We will say that a symmetric rational matrix $Q$ is \emph{positive
semidefinite} when $\Zx^TQ\Zx\ge0$ for every rational vector $\Zx$.
The following argument gives the quadratic inequality and the
principal-submatrix rank bound needed in the core proof.

\begin{lemma}\label{lem:principal}
Suppose $D$ is a $2$-$(15,5,2)$ design, and let $C$ be the matrix
in~\eqref{eq:C-definition}. For a subset $F\sbs\cB$, let $C_F$ be the
principal submatrix indexed by $F$. Then $C_F$ is positive semidefinite
and $\rk C_F\le6$. In particular, if $|F|\ge7$, then $C_F$ is singular.
\end{lemma}
\begin{proof}
Extend a vector on $F$ by zero outside $F$ and apply
\eqref{eq:C-positive}. This proves positive semidefiniteness.
The rank of a submatrix is at most that of the original matrix, so
$\rk C_F\le\rk C\le6$. A square matrix of size $|F|\ge7$ and rank
at most six is singular.
\end{proof}

We have proved the part of Connor's Theorem 3.1 that~\cite{HC} uses
for singularity, with the stronger rank bound. The determinant sign
assertion is recorded separately below. Connor's paper has not been
consulted; these are reconstructions of the statements reported in
Hall and Connor.

\begin{remark}\label{rem:det-nonneg}
For comparison with~\cite{HC}; not used below.
Every symmetric positive semidefinite rational
matrix $Q$ has nonnegative determinant. We prove this by induction on
its size; the determinant of the empty matrix is $1$. The diagonal
entries are nonnegative by testing coordinate vectors. If a diagonal
entry, say $Q(1,1)$, is zero, suppose that $Q(1,i)\ne0$ for some $i$.
For $\tau=-(Q(i,i)+1)/(2Q(1,i))$, the vector
$\tau\Ze_1+\Ze_i$ has quadratic value $-1$, contradicting positive
semidefiniteness. Thus $Q(1,i)=0$ for every $i$, the first row is zero
and $\det Q=0$. In the remaining case $d=Q(1,1)>0$, write
\[
 Q=\pmat{d&\Zu^T\\\Zu&V},\qquad
 Q'=V-d^{-1}\Zu\Zu^T.
\]
Testing $Q$ on $(-d^{-1}\Zu^T\Zy,\Zy)$ gives
$\Zy^TQ'\Zy\ge0$. Subtracting $\Zu(i)/d$ times the first row from
each later row gives $\det Q=d\det Q'$. The induction hypothesis gives
$\det Q'\ge0$, and the result follows. In particular,
$\det C_F\ge0$ for the submatrices in Lemma~\ref{lem:principal}.
\end{remark}

\section{The local restrictions: HC4}

\begin{lemma}\label{lem:claw}
Suppose $D$ is a $2$-$(15,5,2)$ design. Its graph $Z$ has no induced
claw: no block is adjacent to three distinct, pairwise nonadjacent
blocks.
\end{lemma}
\begin{proof}
On a putative centre followed by its three leaves, the principal
submatrix of $C$ would be
\[
 Q_{\mathrm{claw}}=
% BEGIN MATRIX claw
 \begin{pmatrix}
 10&3&3&3\\3&10&-4&-4\\3&-4&10&-4\\3&-4&-4&10
 \end{pmatrix}.
% END MATRIX claw
\]
For $\Zx=(-1,1,1,1)^T$ its quadratic value is
$40+2(-9-12)=-2$, contradicting Lemma~\ref{lem:principal}.
\end{proof}

\begin{remark}\label{rem:source-local}
\cite[Lemmas 4.3 and 4.4, at $k=5$]{HC}
For comparison with Hall--Connor; not used below.
Suppose $D$ is a $2$-$(15,5,2)$ design. If distinct blocks have either
of the following intersection matrices, then $t=1$:
\[
 \pmat{5&1&1&1&1&2\\&5&1&2&2&1\\&&5&2&2&1\\
 &&&5&t&2\\&&&&5&2\\&&&&&5},\qquad
 \pmat{5&1&1&2&2\\&5&1&1&1\\&&5&1&1\\&&&5&t\\&&&&5}.
\]
The omitted entries are determined by symmetry, and $t$ is the
intersection size of the fourth and fifth blocks.

By Lemma~\ref{lem:intersections}, $t\in\{1,2\}$. If $t=2$ in the
first matrix, the first block is the centre of a claw with leaves the
second, fourth and fifth blocks. If $t=2$ in the second matrix, the
second block is the centre of a claw with leaves the first, fourth and
fifth blocks. Both contradict Lemma~\ref{lem:claw}.
\end{remark}

\begin{lemma}\label{lem:five}
\textup{\cite[Lemma 4.5, at $k=5$]{HC}}
Suppose $D$ is a $2$-$(15,5,2)$ design, and five distinct blocks
$B_1,B_2,B_3,Y,W$ have intersection matrix
\[
 \pmat{5&1&1&1&2\\&5&1&2&1\\&&5&2&1\\&&&5&t\\&&&&5},
\]
with the lower triangle supplied by symmetry. Then $t=2$.
\end{lemma}
\begin{proof}
If $t=1$, the principal submatrix of $C$ is
\[
 Q_{\mathrm{five}}=
% BEGIN MATRIX five
 \begin{pmatrix}
 10&3&3&3&-4\\3&10&3&-4&3\\3&3&10&-4&3\\
 3&-4&-4&10&3\\-4&3&3&3&10
 \end{pmatrix}.
% END MATRIX five
\]
For $\Zx=(-1,1,1,1,-1)^T$, the diagonal contribution is $50$ and
the sum over the upper off-diagonal triangle is $-27$.
Thus $\Zx^TQ_{\mathrm{five}}\Zx=50-54=-4$, contradicting
Lemma~\ref{lem:principal}. Lemma~\ref{lem:intersections} leaves only
$t=2$.
\end{proof}

We will also need two seven-by-seven matrices. We give integer
certificates for their inverses so that the argument does not depend
on a determinant expansion or on Connor's computational Lemma 3.1.

\begin{lemma}\label{lem:inverse}
The following rational matrices $Q_1,Q_2$ are nonsingular:
\[
 Q_1=
% BEGIN MATRIX Q1
 \begin{pmatrix}
 10&3&3&3&-4&-4&-4\\
 3&10&3&-4&3&3&3\\
 3&3&10&-4&3&3&3\\
 3&-4&-4&10&-4&-4&-4\\
 -4&3&3&-4&10&3&3\\
 -4&3&3&-4&3&10&3\\
 -4&3&3&-4&3&3&10
 \end{pmatrix},
% END MATRIX Q1
 \qquad
 Q_2=
% BEGIN MATRIX Q2
 \begin{pmatrix}
 10&3&3&3&3&-4&-4\\
 3&10&3&-4&-4&3&3\\
 3&3&10&-4&-4&3&3\\
 3&-4&-4&10&3&-4&-4\\
 3&-4&-4&3&10&-4&-4\\
 -4&3&3&-4&-4&10&3\\
 -4&3&3&-4&-4&3&10
 \end{pmatrix}.
% END MATRIX Q2
\]
\end{lemma}
\begin{proof}
Put
\begingroup
\setlength{\arraycolsep}{4pt}
\[
 K_1=
% BEGIN MATRIX K1
 \begin{pmatrix}
 20&-12&-12&-6&8&8&8\\
 -12&9&7&4&-5&-5&-5\\
 -12&7&9&4&-5&-5&-5\\
 -6&4&4&4&-2&-2&-2\\
 8&-5&-5&-2&5&3&3\\
 8&-5&-5&-2&3&5&3\\
 8&-5&-5&-2&3&3&5
 \end{pmatrix},
% END MATRIX K1
 \qquad
 K_2=
% BEGIN MATRIX K2
 \begin{pmatrix}
 17&-11&-11&-7&-7&6&6\\
 -11&9&7&5&5&-4&-4\\
 -11&7&9&5&5&-4&-4\\
 -7&5&5&5&3&-2&-2\\
 -7&5&5&3&5&-2&-2\\
 6&-4&-4&-2&-2&4&2\\
 6&-4&-4&-2&-2&2&4
 \end{pmatrix}.
% END MATRIX K2
\]
\endgroup
Integer multiplication gives
\begin{equation}\label{eq:inverse-cert}
 K_1Q_1=14I_7,\qquad K_2Q_2=14I_7.
\end{equation}
These are explicit certificates: for each $u,w\in\{1,\ldots,7\}$,
the sum $\sum_{i=1}^7K_j(u,i)Q_j(i,w)$ is $14$ when $u=w$ and $0$
otherwise. For example, the first two entries of $K_1Q_1$ are
\[
 200-36-36-18-32-32-32=14,\qquad
 60-120-36+24+24+24+24=0.
\]
All entries use this same displayed multiplication; the check in
Section~\ref{sec:checks} reads the matrices directly from this source
and checks each entry with integer arithmetic. If $Q_j\Zx=0$, then
$14\Zx=K_jQ_j\Zx=0$. Since $14\ne0$ in $\rats$, $\Zx=0$.
Thus $Q_j$ is nonsingular.
\end{proof}

Both determinants are $235298=2\cdot7^6>0$. We will use their
nonsingularity against the rank bound, not the sign of these
determinants.

\section{The six-cliques: the local part of HC5}

\begin{lemma}\label{lem:local-split}
Suppose $D$ is a $2$-$(15,5,2)$ design and $B_1\sim B_2$ in its block
graph $Z$. There are six-cliques $S_1,S_2$ such that
\[
 B_1,B_2\in S_1,\qquad
 S_1\cap S_2=\{B_1\},\qquad
 S_1\cup S_2=\{B_1\}\cup\Gamma(B_1).
\]
Moreover, $S_1$ is the unique six-clique containing $B_1$ and $B_2$.
\end{lemma}
\begin{proof}
We will partition the blocks other than $B_1,B_2$ into
\[
\begin{aligned}
 \cM&=\Gamma(B_1)\cap\Gamma(B_2),\\
 \cU&=\Gamma(B_1)\diff(\Gamma(B_2)\cup\{B_2\}),\\
 \cV&=\Gamma(B_2)\diff(\Gamma(B_1)\cup\{B_1\}),\\
 \cR&=\cB\diff(\{B_1,B_2\}\cup\cM\cup\cU\cup\cV).
\end{aligned}
\]
By Lemma~\ref{lem:common}, these disjoint sets have sizes $5,4,4,6$,
respectively. In particular,
$\Gamma(B_1)=\{B_2\}\cup\cM\cup\cU$ and
$\Gamma(B_2)=\{B_1\}\cup\cM\cup\cV$, as disjoint unions.

Both $\cU$ and $\cV$ are cliques. For distinct nonadjacent $Y,Y'\in\cU$,
the block $B_1$ would be the centre of a claw with leaves $B_2,Y,Y'$.
For distinct nonadjacent $W,W'\in\cV$, the centre $B_2$ and leaves
$B_1,W,W'$ give the same contradiction to Lemma~\ref{lem:claw}.

Fix $Y\in\cU$. The adjacent pair $B_1,Y$ has five common neighbours.
Inside $\Gamma(B_1)$, the three elements of $\cU\diff\{Y\}$ are
adjacent to $Y$, and $B_2$ is not. Thus exactly two elements of $\cM$
are adjacent to $Y$. There are consequently $4\cdot2=8$ edges between
$\cM$ and $\cU$. Since $|\cM|=5$, some $B_3\in\cM$ has at least
two neighbours in $\cU$; otherwise there would be at most five such
edges.

Put $j=|\cU\diff\Gamma(B_3)|$, so $0\le j\le2$. Let $h,e,\ell$ count
the neighbours of $B_3$ in $\cM\diff\{B_3\}$, $\cV$ and $\cR$,
respectively. Counting common neighbours of $B_1,B_3$, then of
$B_2,B_3$, and finally all neighbours of $B_3$, gives
\[
\begin{aligned}
 1+h+(4-j)&=5,\\
 1+h+e&=5,\\
 2+h+(4-j)+e+\ell&=10.
\end{aligned}
\]
The first $1$ counts $B_2$, the second $1$ counts $B_1$, and the $2$
counts both. Solving gives $h=j$, $e=4-j$ and $\ell=j$.

We will now prove $j=0$. Define
\[
 \cU_0=\cU\diff\Gamma(B_3),\qquad
 \cV_1=\cV\cap\Gamma(B_3).
\]
They have sizes $j$ and $4-j$, and each is a clique. There are no edges
between them. Indeed, for $Y\in\cU_0$ and $W\in\cV_1$, the five
distinct blocks $(B_1,B_2,B_3,Y,W)$ have precisely the hypotheses of
Lemma~\ref{lem:five}: the first three form a triangle; $Y$ is adjacent
to $B_1$ and to neither $B_2$ nor $B_3$; and $W$ is adjacent to $B_2,B_3$
and not to $B_1$. That lemma gives $Y\not\sim W$.

The set
\[
 \Omega=\{B_1,B_2,B_3\}\cup\cU_0\cup\cV_1
\]
is a disjoint union of size $3+j+(4-j)=7$. Order its elements by listing
$B_1,B_2,B_3$, then $\cU_0$, then $\cV_1$. If $j=1$, its principal
matrix $C_\Omega$ is $Q_1$ of Lemma~\ref{lem:inverse}; if $j=2$, it is $Q_2$.
To identify every entry, the first three form a triangle, every member
of $\cU_0$ is adjacent only to $B_1$ amongst these three, every member of
$\cV_1$ is adjacent to $B_2,B_3$ amongst these three, the two latter sets
are separately cliques, and there are no edges between them.
Equation~\eqref{eq:C-entries} gives the displayed matrices. Both are
nonsingular by Lemma~\ref{lem:inverse}, whereas
Lemma~\ref{lem:principal} makes every seven-block principal matrix
singular. Thus $j\notin\{1,2\}$, and $j=0$.

It follows that $B_3$ has no neighbours in $\cM\diff\{B_3\}$ and is
adjacent to all four members of $\cU$. The set
$\cM\diff\{B_3\}$ is a clique. Otherwise two nonadjacent members,
together with $B_3$, would be three pairwise nonadjacent neighbours of
$B_1$. Thus the sets
\[
 S_1=\{B_1,B_2\}\cup(\cM\diff\{B_3\}),\qquad
 S_2=\{B_1,B_3\}\cup\cU
\]
are six-cliques. The disjoint partition above gives their stated
intersection and union.

To prove uniqueness, let $K$ be any six-clique containing $B_1,B_2$.
Its four remaining members lie in $\cM$. It cannot contain $B_3$,
since $B_3$ has no neighbour in $\cM\diff\{B_3\}$, whereas $K$ would
have three further members there. Thus the four remaining members are
exactly $\cM\diff\{B_3\}$, and $K=S_1$.
\end{proof}

\begin{corollary}\label{cor:two-cliques}
Suppose $D$ is a $2$-$(15,5,2)$ design. Every adjacent pair of blocks in
$Z$ belongs to a unique six-clique. Every block belongs to exactly two
six-cliques, whose intersection is that block.
\end{corollary}
\begin{proof}
Lemma~\ref{lem:local-split} applies to every ordered adjacent pair, so
it proves the first assertion. Fix $B_1$. It has ten neighbours, so we
may choose $B_2\in\Gamma(B_1)$ and obtain $S_1,S_2$ from that lemma.
Let $K$ be a six-clique containing $B_1$, and choose
$X\in K\diff\{B_1\}$. The union assertion puts $X$ in exactly one of
$S_1\diff\{B_1\}$ and $S_2\diff\{B_1\}$. Both $K$ and this $S_i$
contain the adjacent pair $B_1,X$, so uniqueness for that pair gives
$K=S_i$. Finally, $S_1$ and $S_2$ are distinct since their intersection
has size one, and both have size six.
\end{proof}

This proves the uniqueness assertion of Hall--Connor Lemma 5.1 at
$k=5$. The explicit use of uniqueness for every adjacent pair is what
lets the local constructions define a single global family.

\section{The global family and the embedding: HC5}

\begin{lemma}\label{lem:global}
Suppose $D$ is a $2$-$(15,5,2)$ design, with block graph $Z$, and let
$\cS$ be the family of all six-cliques of $Z$. Then $|\cS|=7$;
every block belongs to exactly two members of $\cS$; and two distinct
members of $\cS$ have exactly one block in common. Further, for each
$S_0\in\cS$, every point of $D$ occurs in exactly two blocks of $S_0$.
\end{lemma}
\begin{proof}
By definition every member has six blocks, and
Corollary~\ref{cor:two-cliques} gives exactly two members on every
block. Counting incident pairs $(B,S_0)$ gives
$6|\cS|=2|\cB|=42$, whence $|\cS|=7$.

Two different six-cliques cannot share two different blocks: those
blocks would be adjacent and would lie in two six-cliques, contrary to
Corollary~\ref{cor:two-cliques}. To show that they always share one,
fix $S_0\in\cS$. For each $B\in S_0$, let $f(B)$ be the other
six-clique containing $B$. This is well-defined and lies in
$\cS\diff\{S_0\}$. If $f(B)=f(B')$ for different $B,B'\in S_0$,
then $S_0$ and $f(B)$ share these two blocks, which is impossible.
Thus $f$ is injective. Its domain and codomain both have size six, so
it is surjective. Every other six-clique therefore shares a block
with $S_0$, and it shares exactly one.

Fix $S_0\in\cS$ and put
$m(p)=|\{B\in S_0:p\in B\}|$ for $p\in P$. Counting point-block
incidences in $S_0$, and then ordered pairs of distinct blocks of
$S_0$ with a point in their intersection, gives
\[
 \sum_{p\in P}m(p)=6\cdot5=30,\qquad
 \sum_{p\in P}m(p)(m(p)-1)=6\cdot5=30.
\]
The second equality uses that each such pair of blocks meets in one
point. Thus, calculating in $\ints$,
\[
\begin{aligned}
 \sum_{p\in P}(m(p)-2)^2
 &=\sum_{p\in P}m(p)(m(p)-1)-3\sum_{p\in P}m(p)+4|P|\\
 &=30-90+60=0.
\end{aligned}
\]
Each summand is nonnegative, so each is zero. Thus $m(p)=2$ for every
point $p$.
\end{proof}

\begin{theorem}\label{thm:embedding}
\textup{\cite[Theorems 2.1 and 3.2, at $k=5$]{HC}}
Suppose $D$ is a $2$-$(15,5,2)$ design. There is a symmetric
$2$-$(22,7,2)$ design having $D$ as the residual design of one block.
In particular, there is an integer $22\times22$ matrix $M$ with
$MM^T=5I_{22}+2J_{22}$.
\end{theorem}
\begin{proof}
Let $\cS$ be the seven six-cliques from Lemma~\ref{lem:global}.
For each $S_0\in\cS$, adjoin a new point $x_{S_0}$; these seven points
are distinct from one another and from every old point. Replace an old
block $B$ by
\[
 B^*=B\cup\{x_{S_0}:S_0\in\cS,\ B\in S_0\},
\]
and adjoin the new block $B_0=\{x_{S_0}:S_0\in\cS\}$.
There are $15+7=22$ points and $21+1=22$ block labels. Each old block
gets exactly two new points, so every block has size seven. Each old
point still occurs in seven blocks. Each new point occurs in the six
old blocks of its clique and in $B_0$, also seven times.

Two distinct old points occur together twice, as before. An old point
$p$ and a new point $x_{S_0}$ occur together in exactly the two blocks
of $S_0$ containing $p$, by Lemma~\ref{lem:global}. Two distinct new
points $x_{S_0},x_{S_1}$ occur together in the unique old block in
$S_0\cap S_1$ and in $B_0$, exactly twice. These are all types of
pairs, so the construction is a symmetric $2$-$(22,7,2)$ design.
Deleting $B_0$ and its seven points recovers each original labelled
block and its incidences.

We can also record the matrix identity directly. Let $T$ be the
$21\times7$ incidence matrix between $\cB$ and $\cS$,
$T(B,S_0)=1$ exactly when $B\in S_0$. The identities just proved are
\[
 NT=2J_{P,\cS},\qquad
 T^TT=5I_{\cS}+J_{\cS}.
\]
The diagonal entries of the second product are six; its off-diagonal
entries are one. With the old points first, the new points second,
and the new block last, the extended incidence matrix and its product
are
\begin{equation}\label{eq:extended}
\begin{aligned}
 N^*&=\pmat{N&0\\T^T&\Zj_{\cS}},\\
 N^*N^{*T}
 &=\pmat{NN^T&NT\\T^TN^T&T^TT+J_{\cS}}\\
 &=\pmat{5I_P+2J_P&2J_{P,\cS}\\
 2J_{\cS,P}&5I_{\cS}+2J_{\cS}}=5I+2J.
\end{aligned}
\end{equation}
Choose arbitrary enumerations of the $22$ points and $22$ block
labels. Reindexing $N^*$ gives the required integer matrix $M$.
\end{proof}

This is the sufficient direction of Hall--Connor Theorem 2.1 and the
existence part of their Theorem 3.2, specialized to $k=5$. We do not
need uniqueness of the embedding or the converse of Theorem 2.1.

\section{The even-order obstruction: HC6}\label{sec:BRC}

\begin{lemma}\label{lem:det-rank-one}
For $n\ge1$ and rational numbers $\alpha,\beta$,
\[
 \det(\alpha I_n+\beta J_n)=\alpha^{n-1}(\alpha+n\beta).
\]
\end{lemma}
\begin{proof}
Column $i$ is $\alpha\Ze_i+\beta\Zj$. Expand the determinant by multilinearity
in the columns. A term in which two or more columns are chosen from
$\beta\Zj$ is zero, since two columns are proportional. The term using
only the $\alpha\Ze_i$ columns is $\alpha^n$. For each $i$, the term choosing
$\beta\Zj$ in column $i$ and the standard columns elsewhere is
$\beta\alpha^{n-1}$. To see its sign and coefficient, expand
$\Zj=\sum_{j=1}^n\Ze_j$ in that column. All terms except $j=i$ have
two equal standard columns, and the remaining matrix is $I_n$.
There are $n$ choices of $i$, so the determinant is
$\alpha^n+n\beta\alpha^{n-1}=\alpha^{n-1}(\alpha+n\beta)$.
\end{proof}

\begin{lemma}\label{lem:odd-power}
There is no integer $d$ with $d^2=49\cdot5^{21}$.
\end{lemma}
\begin{proof}
We will first prove that $49\cdot5^{2m+1}$ is not an integer square
for every nonnegative integer $m$. Squares modulo five are $0,1,4$,
as follows by squaring the residues $0,1,2,3,4$. In particular, if
$5$ divides $d^2$, then $5$ divides $d$.

For $m=0$, a proposed equality $d^2=49\cdot5$ forces $d=5e$ for an
integer $e$. Cancelling a factor of five gives $5e^2=49$, impossible
modulo five. For the induction step, suppose the assertion is known
for $m$ and that $d^2=49\cdot5^{2m+3}$. Again $d=5e$, and cancelling
$25$ gives $e^2=49\cdot5^{2m+1}$, contradicting the induction
hypothesis. Taking $m=10$ proves the result.
\end{proof}

\begin{theorem}[The needed Bruck--Ryser--Chowla instance]
\label{thm:BRC}
There is no integer $22\times22$ matrix $M$ satisfying
$MM^T=5I_{22}+2J_{22}$.
\end{theorem}
\begin{proof}
Take determinants, viewing the equality also over $\rats$. Determinant
multiplicativity, invariance under transpose, and
Lemma~\ref{lem:det-rank-one} give
\[
 \det(M)^2=\det(MM^T)=\det(5I_{22}+2J_{22})
   =5^{21}(5+22\cdot2)=49\cdot5^{21}.
\]
Since $M$ has integer entries, the permutation formula for its
determinant shows that $\det(M)$ is an integer. This contradicts
Lemma~\ref{lem:odd-power}.
\end{proof}

\begin{proof}[Proof of Theorem~\ref{thm:main}]
If a $2$-$(15,5,2)$ design existed, Theorem~\ref{thm:embedding} would
give the integer matrix excluded by Theorem~\ref{thm:BRC}.
Lemma~\ref{lem:parameters} establishes the equivalence with the five
parameters in the statement.
\end{proof}

Theorem~\ref{thm:BRC} is the even-order case of the obstruction
usually included in the Bruck--Ryser--Chowla theorem. The proof uses
neither Hasse--Minkowski nor a theorem about sums of four squares.
For a symmetric $2$-$(v,k,\lambda)$ design the same determinant
calculation gives $k^2(k-\lambda)^{v-1}$, since
$k+\lambda(v-1)=k^2$. When $v$ is even, the odd exponent $v-1$
forces $k-\lambda$ to be a square by prime factorization. Here we
proved the necessary prime-factor obstruction by an elementary
descent at five. A quadratic-form proof of the odd-order case would
address a different part of the general theorem.

\begin{corollary}\label{cor:complement}
There is no design with $(v,k,b,r,\lambda)=(15,10,21,14,9)$.
\end{corollary}
\begin{proof}
Complement each labelled block in the $15$-point set. Each complement
has size $15-10=5$, and each point belongs to $21-14=7$ complemented
blocks. For distinct points $p,q$, the number of original blocks
containing neither point is, by inclusion--exclusion,
$21-14-14+9=2$. Thus the complement would be a $2$-$(15,5,2)$ design,
contrary to Theorem~\ref{thm:main}.
\end{proof}

\section{Sources, outside facts and checks}\label{sec:checks}

The proof above is complete as an ordinary proof. The following table
records the source of its design-theoretic steps; a source attribution
does not add an unproved hypothesis.

\begin{center}
\begin{tabular}{L{.20\textwidth}L{.32\textwidth}L{.38\textwidth}}
\toprule
Step here & Source & What we prove here\\
\midrule
HC1, Lemma~\ref{lem:intersections} &
Hall--Connor Lemma 4.1, quoting Connor Lemma 4.1 &
All intersection sizes and both multiplicities, by moments.\\[3pt]
HC2, Lemma~\ref{lem:common} &
Hall--Connor Lemma 4.2, quoting Connor Lemma 4.2 &
Both rows of the table, by $G^2=5G+50J$.\\[3pt]
HC3, Lemmas~\ref{lem:Connor}--\ref{lem:principal} &
Connor Theorem 3.1, as used by Hall--Connor &
Nonnegative quadratic forms and rank at most six, including singularity
above size six.\\[3pt]
HC4, Lemmas~\ref{lem:claw} and~\ref{lem:five} &
Hall--Connor Lemmas 4.3--4.5; the last two quote Connor &
The claw and five-block restrictions; the first subsumes the side
patterns in Remark~\ref{rem:source-local}.\\[3pt]
HC4, Lemma~\ref{lem:inverse} &
Hall--Connor (5.3), calculated there using Connor Lemma 3.1 &
The only two nonsingularity cases needed, by integer left inverses.\\[3pt]
HC5, Sections 5--6 &
Hall--Connor Lemma 5.1, Section 6, Theorems 2.1 and 3.2 &
Local uniqueness, the global family, and the embedding at $k=5$.\\[3pt]
HC6, Section~\ref{sec:BRC} &
Bruck--Ryser--Chowla, even-order case &
The determinant formula and the obstruction at five, in full.\\
\bottomrule
\end{tabular}
\end{center}

We did not consult Connor's paper; we use only the statements
reported in Hall--Connor and prove all their required $k=5$ instances.
We do not need Connor's general determinant tool, the $k=6$
nonexistence result, or the examples cited in the introduction of
Hall--Connor. In~\cite{HC}, Lemmas 4.1
and 4.2 are on page 37; the local matrices (4.1), (4.4) and (4.5)
are on pages 38--39. Equation (5.3) is printed on page 40 as
\[
 |C_{k+2}|=j(j-k+2)(k-6)(k+2)^{k+1}.
\]
At $k=5$, this is $j(3-j)7^6$, hence $2\cdot7^6$ for both $j=1$
and $j=2$.
The even-order determinant argument also appears in
\cite[Theorem 8.3.1]{BvM}. Its proof
is included in Section~\ref{sec:BRC} above.

There are two details worth recording for the structured proof.
Hall--Connor's count of entries in their array $C^*$ occurs before
they have proved the relevant set is a clique. Our proof establishes
that clique property first, then obtains exactly eight edges.
Also, we explicitly define $\cS$ to be \emph{all} six-cliques and
prove the adjacent-pair uniqueness that identifies the local choices.
The two seven-block cases cannot be excluded by nonnegative
principal minors alone; the rank assertion in HC3 is necessary.

The outside mathematical facts used without proof are elementary
finite counting, the pigeonhole principle, matrix rank under transpose
and multiplication by the transpose over $\rats$, the rank inequality
for a zero product, rank monotonicity for submatrices, the full-rank
criterion for invertibility, and basic determinant properties:
multilinearity, alternation, multiplicativity and invariance under
transpose. Nonnegativity of rational squares,
integer division with remainder, and induction supply the arithmetic.
We prove the special determinant formula and the prime obstruction
in the core text. The nonnegative-minor statement is the side result
in Remark~\ref{rem:det-nonneg}; neither that remark nor
Remark~\ref{rem:source-local} is a dependency of the core proof or
of its structured and Lean versions. We use no
classification of strongly regular graphs, no design existence table,
and no unproved theorem of Connor or Hall--Connor.

\subsection*{Exact numerical checks}

The executable check is \code{checks/check_proof.py}, run from this
directory by
\begin{lstlisting}
python checks/check_proof.py
\end{lstlisting}
for a Python 3 with SymPy. We used Python 3.12.14 and SymPy 1.14.0. All
computations use exact integers or rational numbers, and the script
prints its results. The check reads the six labelled certificate
matrices directly from \code{hallconnor.tex}, checks the two quadratic
values $-2,-4$, and checks all $49$ entries of each inverse identity
in~\eqref{eq:inverse-cert}. It constructs the two seven-block patterns
independently from their adjacency rules and compares every entry.
It also verifies their determinants, their positive leading principal
minors, and the rank-six boundary case $j=0$.

Further checks cover the intersection counts, the algebra giving
$C^2=35C$ and $CN^T=0$, and the determinant
\[
 \det(5I_{22}+2J_{22})=23365020751953125=49\cdot5^{21}.
\]
An integer square-root check and division by five independently
confirm the arithmetic obstruction. The script also calls
\code{checks/check_t7.py}, which constructs the triangular graph
$T(7)$ directly from the two-subsets of a seven-set. It checks both
rows of the common-neighbour table, absence of claws and every
five-block pattern. For all $210$ ordered adjacent pairs, it checks
each step of Lemma~\ref{lem:local-split}: the partition sizes
$5,4,4,6$, the eight edges, the choice of $B_3$, the counts
$(h,e,\ell)=(j,4-j,j)$, the two six-cliques, their intersection and
union, and uniqueness. It also enumerates all six-cliques and checks
their incidence identities. In every local construction $j=0$;
the excluded $j=1,2$ matrices are checked by their separate certificates.

The biplane controls repeat the residual and reconstruction calculations
for the $(11,5,2)$ and $(16,6,2)$
biplanes. Their residual block sizes are three and four, so they do
not test the structural argument at block size five. No
$2$-$(15,5,2)$ design exists to test: $T(7)$ supplies a graph control but
no integral incidence matrix $N$ for the hypothetical design. The
formula checks in \code{check_counts_and_algebra} re-evaluate the
displayed algebra; they do not independently prove the counts. The
controls cannot establish the quantified structural argument or the
design-to-matrix bridge. The contradiction requires the integrality
of $N$, beyond the consistent graph structure of $T(7)$.

These checks do not search over designs and do not replace any
quantified argument in the proof. The largest individual certificate
matrices are $7\times7$; the $22\times22$ determinant is a diagnostic
check only. In Lean its value is proved symbolically, as in
Lemma~\ref{lem:det-rank-one}, not by expanding $22!$ terms.

\subsection*{What Mathlib contains}

We inspected Mathlib at commit
\code{082e2d37e8b0463410cdb532e111cd43d5a66174} (4 September 2026),
the version pinned in \code{lake-manifest.json}, with Lean
\code{v4.34.0-rc2}. Paths here are relative to
\code{Mathlib/}.

\begingroup
\raggedright
\begin{itemize}
\item \code{Matrix.det_mul}, \code{Matrix.det_transpose} and
\code{Matrix.det_smul}, together with \code{Matrix.det_reindex_self}, are in
\code{LinearAlgebra/Matrix/Determinant/Basic.lean}.
The determinant lemma
\code{Matrix.det_one_add_replicateCol_mul_replicateRow} is in
\code{LinearAlgebra/Matrix/SchurComplement.lean}, line 418.
It states that the determinant of $I+\Zu\Zv^T$ is
$1+\Zv^T\Zu$. Scaling $5I+2J=5(I+(2/5)J)$ over $\rats$ gives the
needed determinant calculation.
The more general \code{Matrix.det_add_replicateCol_mul_replicateRow}
is in the same file, line 428.
\item In \code{LinearAlgebra/Matrix/Rank.lean},
\code{Matrix.rank_add_rank_le_card_of_mul_eq_zero} (line 576) gives
$\rk C+\rk N^T\le21$ from $CN^T=0$.
\code{Matrix.rank_self_mul_transpose} (line 561) applies over the
linearly ordered field $\rats$ and gives $\rk(NN^T)=\rk N$;
\code{Matrix.rank_transpose} gives $\rk N^T=\rk N$.
\code{Matrix.rank_of_isUnit} (line 211), with
\code{Matrix.isUnit_iff_isUnit_det} from
\code{LinearAlgebra/Matrix/NonsingularInverse.lean}, turns
$\det(NN^T)=5^{14}\cdot35\ne0$ into $\rk N=15$.
Thus Lemma~\ref{lem:det-rank-one} at $n=15$ and $n=22$ serves both
rank and obstruction. The file also contains
\code{Matrix.rank_submatrix_le} and \code{Matrix.rank_mul_le}.
No separate injectivity or \code{LinearMap} argument is needed.
The singularity equivalence \code{Matrix.exists_mulVec_eq_zero_iff}
is in \code{LinearAlgebra/Matrix/ToLinearEquiv.lean}, line 163.
\item The rational quadratic-form proof can use sums of squares
directly. Mathlib also has
\code{Matrix.posSemidef_iff_dotProduct_mulVec} and
\code{Matrix.PosSemidef.submatrix} in
\code{LinearAlgebra/Matrix/PosDef.lean}.
\code{Matrix.PosSemidef.det_nonneg}, in
\code{Analysis/Matrix/PosDef.lean}, line 48, assumes an
\code{RCLike} scalar field; it is not directly a theorem over
$\rats$. The elementary elimination proof is retained only in
Remark~\ref{rem:det-nonneg}; it is outside the core Lean chain.
The quadratic inequality and rank bound alone suffice for that chain.
\item \code{padicValNat.mul} and \code{padicValNat.pow} are in
\code{NumberTheory/Padics/PadicVal/Basic.lean}, lines 424 and 441.
They give an alternative to our descent: take the natural absolute
value of $\det M$, install \code{Fact (Nat.Prime 5)}, establish that
the factors are nonzero, and deduce $2\nu_5(|\det M|)=21$.
\item In \code{NumberTheory/SumFourSquares.lean}, line 199, Mathlib has
the precise statement
\begin{lstlisting}
Nat.sum_four_squares (n : Nat) :
  Exists fun a : Nat => Exists fun b : Nat =>
  Exists fun c : Nat => Exists fun d : Nat =>
    a^2 + b^2 + c^2 + d^2 = n
\end{lstlisting}
Here we spell the existential quantifiers with \code{Exists}.
The multiplicative identity
\code{Nat.euler_four_squares} is also present.
Neither declaration is a dependency of this proof. A search of the
\code{Mathlib/} source found no declaration of
Hasse--Minkowski or Bruck--Ryser--Chowla under those names. No such
declaration is needed.
\end{itemize}
\endgroup

\section{The Lean statements}\label{sec:Lean}

We will use the structure name \code{Design}, as required by
\code{cgss/notation.md}, Sections 7 and 10. This section gives the
interface and the signatures that the Lean formalization, in
\code{LineGraphs/Cited/HallConnor/}, implements, together with the plan
of their proofs that we wrote before the formalization.

We will keep points and block labels as separate finite types, so
that repeated blocks are allowed by the interface. One suitable
shape is
\begin{lstlisting}
Design (P B : Type*) [Fintype P] [Fintype B]
    [DecidableEq P] [DecidableEq B] (v k b r lam : Nat)
\end{lstlisting}
Here all five numbers \code{v k b r lam} are explicit parameters.
The fields \code{card_points : Fintype.card P = v} and
\code{card_blocks : Fintype.card B = b} link the finite types to the
parameters. There is a field \code{block : B -> Finset P}, a proof that every block
has size \code{k}, a proof that every point occurs in \code{r}
blocks, and a proof that each distinct pair occurs in \code{lam}
blocks. It must have no injectivity field for \code{block}.
We will use \code{D.incMatrix} for the integer incidence matrix,
\code{D.compl} for the complement and \code{D.IsSymmetric} for the
condition \code{v = b}, exactly as in the notation table.
Lemma~\ref{lem:parameters} converts an unbundled
$2$-$(15,5,2)$ hypothesis, without a prescribed replication number or
block count, to this shape. An arbitrary $21$-element label type can
then be reindexed by \code{Fin 21}.

The following are the declaration signatures. Numeric suffixes name parameter values, not
theorem numbers from the paper. We use \code{False} for a contradiction
from a supplied design. The embedding signature returns the integer
matrix and equation~\eqref{eq:extended} that the obstruction consumes.
\begin{lstlisting}
theorem Design.not_nonempty_15_5_2 :
  Not (Nonempty (Design (Fin 15) (Fin 21) 15 5 21 7 2))

theorem Design.false_of_parameters_15_5_2
    (D : Design P B 15 5 21 7 2) : False

theorem Design.hall_connor
    (D : Design P B 15 5 21 7 2) :
    Exists fun M : Matrix (Fin 22) (Fin 22) Int =>
      M * M.transpose =
        5 * (1 : Matrix (Fin 22) (Fin 22) Int) +
        2 * Matrix.of (fun _ _ => (1 : Int))

theorem Matrix.bruck_ryser_chowla_twenty_two
    (M : Matrix (Fin 22) (Fin 22) Int) :
    Not (M * M.transpose =
      5 * (1 : Matrix (Fin 22) (Fin 22) Int) +
      2 * Matrix.of (fun _ _ => (1 : Int)))
\end{lstlisting}
In the matrix signatures, the scalar-looking factors
\code{5} and \code{2} here are scalar identity matrices under the
matrix numeral convention, so their matrix products give the intended
scalar multiples. Using \code{5} and \code{2} as scalar actions instead
is equally suitable. For the abstract signature, the finite and
decidable-equality instances for \code{P} and \code{B} are implicit
section variables. The Lean formalization uses the scalar actions.

The conclusion of \code{Design.hall_connor} is precisely the matrix
hypothesis of \code{Matrix.bruck_ryser_chowla_twenty_two}.

We will retain the residual assertion with
\code{Design.hallConnorExtension}. It has point type
\code{Sum P D.sixCliques} and block type \code{Option B}, where
the finite set \code{D.sixCliques} is coerced to its subtype of
members, and \code{none} is the new block. Its five parameters are
\code{22 7 22 7 2}. Its incidence relation is exactly the case
distinction in Theorem~\ref{thm:embedding}. Projection onto the
old summands recovers \code{D}; this avoids making the main
nonexistence theorem depend on a general residual-design API.

The internal names will describe their conclusions. Definitions use
lower camel case; theorems use snake case and live in the namespace of
their object. The settled incidence name \code{D.incMatrix} takes
precedence over the paper's letter $N$; the other objects $G,C,Z$
keep their letters. The parameters of a design are
\code{v k b r lam}; \code{lam} avoids Lean's reserved lambda token.
A concrete name map is as follows.

\begin{center}
\begin{tabular}{L{.32\textwidth}L{.58\textwidth}}
\toprule
Object or conclusion & Declaration\\
\midrule
Incidence and intersection matrices &
\code{Design.incMatrix}, \code{Design.G}\\
Connor matrix and block graph &
\code{Design.C}, \code{Design.Z}\\
Intersection restriction &
\code{Design.inter_card_eq_one_or_two}\\
Degree and common neighbours &
\code{Design.card_neighbors}, \code{Design.card_commonNeighbors}\\
Connor identities &
\code{Design.C_mul_C}, \code{Design.C_mul_incMatrix_transpose}\\
Quadratic inequality and rank &
\code{Design.C_quadratic_nonneg}, \code{Design.rank_C_le_six}\\
Seven-block singularity &
\code{Design.det_C_submatrix_eq_zero}\\
Local restrictions &
\code{Design.Z_clawFree}, \code{Design.not_adj_of_five_block_pattern}\\
Clique uniqueness &
\code{Design.exists_unique_sixClique_of_adj}\\
Global family &
\code{Design.sixCliques}, \code{Design.card_sixCliques}\\
Clique incidence counts &
\code{Design.card_sixCliques_through_block},
\code{Design.sixClique_inter_card},
\code{Design.card_blocksThrough_in_sixClique}\\
Finite certificate values &
\code{Computations.C19.claw_quadratic},
\code{Computations.C19.five_block_quadratic},
\code{Computations.C19.left_inverse_Q1},
\code{Computations.C19.left_inverse_Q2}\\
Complement needed by CGSS &
\code{Design.compl}, \code{Design.not_nonempty_15_10_9}\\
\bottomrule
\end{tabular}
\end{center}

Theorems in this table concerning $D$ carry the complete hypotheses
$D:$ \code{Design P B 15 5 21 7 2}, as well as the
blocks or cliques mentioned in their conclusions. The computation
namespace is the C19 namespace already assigned in the project;
the paper's lemma numbers appear in docstrings rather than in
theorem names. The two named external theorems retain their names.

\subsection*{Formalization cost and proof boundaries}

The expensive part will be the finite incidence and clique bookkeeping.
In particular, Lemma~\ref{lem:common} is the densest counting step:
we will split the entry $(G^2)(B_1,B_2)$ into the terms $B=B_1$,
$B=B_2$ and the sum over the other $19$ labels, then express the
three weighted sums over the four filtered finite sets counted by
$x_{ij}$. Along with their total cardinality, these give the four
displayed equations. This will require explicit exclusions, disjointness
and casts from natural cardinalities to integer or rational sums.
The local-to-global clique argument will also need disjoint partitions,
double counts with block labels and exact filtered-set cardinalities.

The matrix rank lemmas listed above make the rank deduction short.
The remaining transport cost is the principal-submatrix bound and
an explicit equivalence from a seven-block subset to \code{Fin 7},
so that the entry pattern and singularity agree with the certificates.

The passage from Theorem~\ref{thm:embedding} to
Theorem~\ref{thm:BRC} has a separate design-to-matrix cost.
If the embedding returns a bundled design, we must apply DS1,
$NN^T=(r-\lambda)I+\lambda J$, to its incidence matrix, and reindex
the point and block types to \code{Fin 22}. A determinant taken on
the original point type must be transported by
\code{Matrix.det_reindex_self}. Our chosen \code{Design.hall_connor}
signature instead returns the matrix identity directly: we will prove
the block product~\eqref{eq:extended} and reindex its rows and columns
to \code{Fin 22}. This still requires compatibility of reindexing with
multiplication, transpose, identity and all-ones matrices; it avoids
making the contradiction depend on a second application of DS1 to a
bundled extension.

The small certificate identities should be inexpensive. We will use
kernel-checked finite matrix multiplication, rational arithmetic,
and linear combinations of rows. The moment polynomials must be
evaluated in $\ints$ or $\rats$, not with truncated subtraction in
$\nn$. The determinant obstruction should also be inexpensive once
the general rank-one determinant identity is instantiated; either the
descent of Lemma~\ref{lem:odd-power} or the existing valuation lemmas
will finish it. We will keep $N^*$ integral and cast it to $\rats$
only for the determinant identity.

There is no unresolved mathematical step in this proof. The structured
proof is \code{cgss/proofs/hallconnor/hallconnor-lamport.tex}. The Lean
formalization, in the folder \code{LineGraphs/Cited/HallConnor/}, implements
the \code{Design} interface and these declarations. Neither needs a
quadratic-form local-global principle or a development of sums of four
squares.

\begin{thebibliography}{9}
\bibitem{BvM}
A.~E. Brouwer and H. Van Maldeghem, \emph{Strongly Regular Graphs},
Cambridge University Press, 2022.
\bibitem{CGSS}
P.~J. Cameron, J.~M. Goethals, J.~J. Seidel and E.~E. Shult,
\emph{Line graphs, root systems, and elliptic geometry},
J. Algebra \textbf{43} (1976), 305--327.
\bibitem{Connor}
W.~S. Connor, Jr., \emph{On the structure of balanced incomplete block
designs}, Ann. Math. Statist. \textbf{23} (1952), 57--71.
We use the statements reproduced in~\cite{HC}, with proofs given above.
\bibitem{HC}
M. Hall, Jr. and W.~S. Connor,
\emph{An embedding theorem for balanced incomplete block designs},
Canad. J. Math. \textbf{6} (1954), 35--41.
\href{https://doi.org/10.4153/CJM-1954-005-4}{doi:10.4153/CJM-1954-005-4}.
\end{thebibliography}
\end{document}
